\chapter{The dot product : \(\vec{A} \cdot \vec{B}\)}


\section{What does the dot product actully mean?}

The dot product operates on two arguments and produces a scalar value as a result. Conceptually, the dot product
measures how much its two arguments point in the same direction. Put another way, the dot product can be thought of as
measuring how much of one of these arguments would be ``projected'' onto the other. The dot product effectively asks,
``how much do the two arguments agree in direction'', or put another way, ``how much do they align''? \\

To help try and illustrate what this means, consider the scenario where we have two arguments. Let us call these two
arguments \(\vec{A}\) and \(\vec{B}\). Furthermore, let us split the argument \(\vec{B}\) into two components relative
to the argument \(\vec{A}\), and let us do it in such a way that one of these components is parallel to argument
\(\vec{A}\), while the other component is perpendicular to argument \(\vec{A}\). In this case, we can define the
parallel component as follows;

\begin{equation*}
    |B|\cos{(\theta)}
\end{equation*}
\vspace{0.0\baselineskip}

This parallel component is a measure of how much of argument \(\vec{B}\) points in the direction of argument \(\vec{A}\).
Alternatively, we can say that it is a measure of how much of argument \(\vec{B}\) is projected onto argument
\(\vec{A}\). If we multiply this parallel component by \(||\vec{A}||\), i.e. the magnitude or length of the argument
\(\vec{A}\), then we get the dot product between \(\vec{A}\) and \(\vec{B}\). \\

\begin{figure}[h]
    \centering
    \includegraphics[width=0.4\textwidth]{images/dot_product_projection.png}
    \caption{Graphical illustration of the parallel component of \(\vec{B}\).}
    \label{fig:cross_product}
\end{figure}
